\documentclass[11pt]{article}
\usepackage[margin=1in]{geometry}
\usepackage{amsmath,amssymb,amsthm,hyperref}
\setlength{\emergencystretch}{2em}
\newtheorem{theorem}{Theorem}
\title{Fixed spectrum, unbounded Krylov exponential error\\Disproof of Conjecture 00000005111}
\author{ziangni-sys}
\date{October 8, 2026}
\begin{document}
\maketitle

\paragraph{The obstruction.}
The source asserts a matrix-exponential Krylov error bound controlled by
the spectral interval and dimension. Neither language version imposes
normality or bounds nonnormality. We hold the spectrum, ambient dimension,
Krylov dimension, time, and starting-vector norm fixed, while the absolute
error of even the best Krylov approximation becomes arbitrarily large.
Thus no finite bound depending only on the named spectral data and
dimensions can hold as written, regardless of its proposed hypergeometric
formula. This does not refute bounds using the numerical range, matrix norm,
eigenvector conditioning, or a normality assumption; those contain additional
information absent from the claim.

\paragraph{Fixed data.}
For $t>0$ take the complex matrix and unit starting vector
\[
 A_t=\begin{pmatrix}0&t\\0&0\end{pmatrix},\qquad
 b=\begin{pmatrix}0\\1\end{pmatrix},\qquad \|b\|_2=1.
\]
The ambient dimension is always two. The first Krylov space, at dimension
one, is the actual space
\[
 \mathcal K_1(A_t,b)=\operatorname{span}_{\mathbb C}\{A_t^0b\}
   =\left\{\begin{pmatrix}0\\c\end{pmatrix}:c\in\mathbb C\right\}.
\]
Since $b\ne0$, its dimension is exactly one. Moreover,
\[
 \det(zI-A_t)=z^2,\qquad \sigma(A_t)=\{0\}
\]
for every $t$. The spectral interval is therefore the fixed interval
$[0,0]$; in particular the spectrum is purely imaginary as well as real.
For $t>0$ these matrices are nonnormal, as
$A_tA_t^*=\operatorname{diag}(t^2,0)$ and
$A_t^*A_t=\operatorname{diag}(0,t^2)$.

\begin{theorem}
At the fixed time $1$, the best approximation to $e^{A_t}b$ from
$\mathcal K_1(A_t,b)$ has Euclidean error $t$. The one-step Arnoldi
approximation attains this error. Consequently these errors have no
finite uniform upper bound with the above data fixed.
\end{theorem}
\begin{proof}
Direct multiplication gives $A_t^2=0$, so every power of order at least
two vanishes. Hence the defining matrix exponential series is a finite
sum, and its convergence and value are explicit:
\[
 e^{A_t}=\sum_{k=0}^{\infty}\frac{A_t^k}{k!}=I+A_t,
       \qquad e^{A_t}b=\begin{pmatrix}t\\1\end{pmatrix}.
\]
For any $v=(0,c)^T\in\mathcal K_1(A_t,b)$,
\[
 \|e^{A_t}b-v\|_2
   =\sqrt{t^2+|1-c|^2}\geq t.
\]
Taking $c=1$ attains equality. This is also the one-step Arnoldi output:
the normalized basis is $V_1=b$, its compressed matrix is
$H_1=b^*A_tb=0$, and $V_1e^{H_1}\|b\|_2=b$.
Finally, for any proposed finite bound $C\in\mathbb R$, choose
$t=|C|+1$. Then every vector in the first Krylov space has error at least
$t>C$, while all the stated fixed data remain unchanged.
\end{proof}

\paragraph{Consequences and scope.}
The first bound clause is false, so its further claims about an optimal
constant, equality cases and a matching construction cannot make the
conjunction true. The example concerns the standard absolute Euclidean
error for approximating $e^A b$, with unit initial vector. A bound allowed
to depend on nonnormality or a field-of-values region is not excluded:
those quantities change with $t$. Likewise a theorem restricted to normal
matrices would have different hypotheses. If ``spectral interval'' were
instead intended to mean a numerical-range enclosure, it would no longer
be an interval determined by the spectrum as stated.

\paragraph{Formal verification.}
The Lean project uses genuine $2\times2$ complex matrices, Mathlib's
algebraic spectrum and canonical \texttt{NormedSpace.exp}, and the
actual complex submodule generated by $b$. The exponential is evaluated
from its defining series using vanishing powers, not assumed as a
certificate. The Euclidean norm is written explicitly as
$\sqrt{|v_0|^2+|v_1|^2}$. The formalization proves the fixed spectrum,
unit input, exponential action, zero compressed Arnoldi matrix, exact
Arnoldi error and a lower bound for \emph{every} vector in the Krylov
space. The final theorem \texttt{no\_spectral\_only\_bound} supplies, for
every real $C$, a matrix with the same spectrum for which all those errors
exceed $C$. No numerical auxiliary calculation is required.
\end{document}
